\chapter{Cybercrime}

Cybercrime can be defined as "The illegal usage of any communication device to commit or facilitate in committing any illegal act". A cybercrime is a type of crime that targets or uses a computer or a group of computers under one network for the purpose of harm.


\section{Classifications of Cybercrimes}
Cybercrimes in general can be classified into four categories:
\begin{enumerate}[itemsep=0.3em, topsep=-0.5em]
   \item \textbf{Individual Cybercrimes:} Target individuals (phishing, spoofing, spam, cyberstalking, etc.).
   \item \textbf{Organization Cybercrimes:} Target organizations (malware attacks, denial of service).
   \item \textbf{Property Cybercrimes:} Target property such as credit cards or intellectual property (digital fraud, copyright breach).
   \item \textbf{Society Cybercrimes:} Target society at large, including cyber-terrorism (most dangerous form).
\end{enumerate}

\section{Cyber Offences under Bangladeshi Laws}
\begin{itemize}[parsep=0.3em]
    \item Hacking or illegal access to critical information infrastructure
    \item Illegal access to computers or computer systems
    \item Destruction of cyberspace or computer systems
    \item Gambling in cyberspace
    \item Digital fraud
    \item Cyber terrorism
    \item Harassment or blackmailing in cyberspace
    \item Spreading hate speech in cyberspace
    \item Aiding in committing cyber offences
\end{itemize}

\section{Digital Signature}

\textbf{Digital signature} means electronic data which:  
\begin{enumerate}[topsep=-0.5em, itemsep=0.5em]
    \item Is related to any other electronic data directly or logically; and
    \item Satisfies the following conditions for validation:
    \begin{enumerate}
        \item Affixed with the signatory uniquely;
        \item Capable of identifying the signatory;
        \item Created securely or under the sole control of the signatory;
        \item Linked to the attached data to detect any alterations.
    \end{enumerate}
\end{enumerate}